\chapter*{Resumen}
Un servicio de correo transaccional suele ofrecer dos interfaces: la API REST y el relé SMTP,
que recibe el correo por el protocolo estándar de envío. Los términos de consumo del servicio
son el precio y la información de cómo, cuándo, cuánto y dónde se consume. Cada interfaz puede
tener límites propios: en el plan Starter de MailerSend, el relé SMTP acepta como máximo cinco
conexiones simultáneas. Para saber qué proveedor admite el envío que necesita una aplicación, hay que
reunir a mano los términos de cada interfaz, y esos términos no se comparan de forma
sistemática.

El \textit{Pricing Intelligence} es la línea de investigación que convierte la información de
los planes de precios en conocimiento para elegir plan o proveedor. En esa línea, dos
Trabajos Fin de Máster del grupo de investigación ISA aportan el \textit{API Capacity Model}
(ACM), que describe cuánto deja consumir cada plan y en qué periodo. Definen también la
Datasheet, un documento legible por máquina con el precio y los límites de una API, escrito en
su esquema YAML, ACM2YAML. Sobre ella, su motor PRIME4API calcula el plan más barato para un
volumen o el tiempo mínimo para enviarlo.

El objetivo de este Trabajo Fin de Máster es llevar ese instrumental a un dominio real y a escala.
Para ello se construye un catálogo de Datasheets de la API REST de ocho servicios de correo
transaccional. Cada dato del catálogo lleva su fuente y su fecha.
El catálogo se ha propuesto como contribución a SPHERE, la plataforma del grupo que publica
planes de precios.

El catálogo lo construye un proceso agéntico diseñado y supervisado por el autor. Unos agentes
de inteligencia artificial redactan la Datasheet de cada API con el precio y los límites que
extraen de las fuentes. Otros agentes revisan
cada Datasheet, y lo que queda abierto lo resuelve el autor.

El envío suele tener dos límites: una cuota (el volumen máximo en un periodo, como un mes) y
un \textit{rate} (las peticiones que admite la API por segundo o por minuto). En seis de las
ocho APIs, con la cuenta aprobada, al ritmo máximo del \textit{rate} la cuota se agota antes
de que acabe su periodo, así que el envío se dimensiona con la cuota. El estado de la cuenta
(en revisión, en pruebas o por su reputación) cambia esos límites en seis APIs, no todas las anteriores,
y ninguna publica el umbral que lo provoca.

El estudio de las ocho APIs muestra lo que ACM2YAML no puede expresar en sus campos, como un
plan que se vende en varios volúmenes, cada uno con su precio, o los topes del relé SMTP. Para
cubrir parte de ello se proponen extensiones del esquema, cada una con el caso real que la motiva.
Este trabajo establece un método de comparación que responde qué APIs admiten un patrón de
envío y cuánto cuesta en cada una, y lo aplica a un caso de ejemplo: una aplicación que envía
100.000 correos al mes.

\vspace{.5cm}
{\raggedright\noindent\textbf{Palabras clave:} Pricing Intelligence, APIs de correo electrónico, \textit{datasheets} de APIs, capacidad de APIs,
\textit{rate limits}, cuotas de uso, comparación sistemática de APIs, ACM2YAML,
propiedades no funcionales, selección de servicios\par}
